\chapter{Wrist Pain}

\section*{Definition}
Pain or discomfort localized to the wrist joint and surrounding structures, commonly resulting from overuse, trauma, or inflammatory conditions.

\section*{Pathophysiology}
The wrist's complex anatomy of eight carpal bones, multiple ligaments, and flexor/extensor tendons makes it susceptible to injury. Pain localizes anatomically: radial side (De Quervain tenosynovitis, scaphoid fracture, thumb arthritis), ulnar side (TFCC tear, extensor carpi ulnaris tendinopathy), dorsal (ganglion cyst, intersection syndrome), volar (carpal tunnel syndrome).

\section*{Etiology}
\begin{itemize}
  \item De Quervain tenosynovitis
  \item Carpal tunnel syndrome
  \item Wrist sprain or fracture (scaphoid, distal radius)
  \item Osteoarthritis or post-traumatic arthritis
  \item Ganglion cyst
  \item Rheumatoid arthritis
  \item Gout or pseudogout
  \item Ligamentous injury (TFCC tear)
  \item Tendinopathy (extensor carpi radialis, flexor carpi ulnaris)
  \item Kienböck disease (avascular necrosis of lunate)
\end{itemize}

\section*{Indications for Evaluation}
Seek care if:
\begin{itemize}
  \item Severe or worsening symptoms
  \item Wrist pain after injury, inability to move the wrist, severe swelling, numbness in fingers, or fever
  \item Accompanied by other concerning signs
\end{itemize}

\section*{Related Symptoms}
\begin{itemize}
  \item Hand numbness
  \item Finger pain
  \item Swelling
  \item Weak grip
  \item Snapping sensation
\end{itemize}
\textbf{Note:} Always consider serious underlying conditions when evaluating persistent wrist pain.

\section*{Self-Care / Home Management}
\begin{itemize}
  \item Rest and avoid activities that may aggravate the condition.
  \item Maintain adequate hydration and a balanced diet.
  \item Over-the-counter remedies may provide symptomatic relief (consult a pharmacist or doctor).
  \item Monitor symptoms and seek professional advice if they persist or worsen.
  \item Avoid self-medicating with prescription medications without medical guidance.
\end{itemize}

\section*{Diagnostic Approach}
\begin{itemize}
  \item A thorough history and physical examination are the first steps in evaluation.
  \item Laboratory tests (blood work, urinalysis) may be ordered based on clinical suspicion.
  \item Imaging studies (X-ray, ultrasound, CT, MRI) may be indicated when structural pathology is suspected.
  \item Specialist referral is arranged when the diagnosis remains uncertain or requires specific expertise.
\end{itemize}

\section*{Treatment / Management Overview}
\begin{itemize}
  \item Treatment targets the underlying cause identified through diagnostic evaluation.
  \item Supportive care and symptom management are provided while awaiting definitive therapy.
  \item Medications, lifestyle modifications, and physical therapies may be prescribed as appropriate.
  \item Follow-up ensures treatment effectiveness and allows timely adjustments to the care plan.
\end{itemize}

\clearpage